\section{Rotating Reference Frames}

Rotating reference frames are important examples of non-inertial frames.

\subsection{Newton's Equations in a Rotating Reference Frame}

We first set up the problems, defining some notations.

\begin{definition}[Notation on Rotating Frame]
    \label{def:rotating-frame-notation}%
    We define:
    \begin{itemize}
        \item \(\vb{e}_1, \vb{e}_2, \vb{e}_3\) as the axes of an inertial frame \(S\);
        \item  \(\vb{e}_1', \vb{e}_2', \vb{e}_3'\) as the axes of a non-inertial rotating frame \(S'\).
    \end{itemize}
    
    From the perspective of the inertial frame, let \(\vb*{\omega}\) be the angular velocity the \(\vb{e}_i'\) axes rotate, so
    \[
        \dot{\vb{e}}_i' = \vb*{\omega} \cross \vb{e}_i'.
    \]

    In the two frames, the position of a particle is, respectively,
    \[
        \vb{x} = x_i \vb{e}_i = x_i' \vb{e}_i'.
    \]

    The differential operators
    \[
        \pqty{\dv{t}}_S, \pqty{\dv{t}}_{S'}
    \]
    means differentiating by Cartesian components with respect to a certain frame.
\end{definition}

\begin{proposition}[Velocity in a Rotating Frame]
    \label{prop:velocity-rotating-frame}%
    The velocity described in a rotating frame is given by
    \[
        \pqty{\dv{\vb{x}}{t}}_{S} = \pqty{\dv{\vb{x}}{t}}_{S'} + \vb*{\omega} \cross \vb{x}.
    \]
\end{proposition}

\begin{proof}
    On one hand, in the inertial frame, we have
    \[
        \dot{\vb{x}} = \dot{x}_i \vb{e}_i = \pqty{\dv{\vb{x}}{t}}_{S},
    \]
    while in the rotating frame, differentiating gives
    \begin{align*}
        \dot{\vb{x}} &= \dot{x}_i' \vb{e}_i' + x_i' \dot{\vb{e}}_i' \\
        &= \dot{x}_i' \vb{e}_i' + x_i' \pqty{\vb*{\omega} \cross \vb{e}_i'} \\
        &= \dot{x}_i' \vb{e}_i' + \vb*{\omega} \cross \vb{x} \\
        &= \pqty{\dv{\vb{x}}{t}}_{S'} + \vb*{\omega} \cross \vb{x}.
    \end{align*}
\end{proof}

This result is expected -- the difference is just the relative velocity of the two frames.

\begin{proposition}[Acceleration in a Rotating Frame]
    \label{prop:acceleration-rotating-frame}%
    The acceleration in a rotating frame is given by
    \[
        \boxed{\pqty{\dv[2]{\vb{x}}{t}}_{S} = \pqty{\dv[2]{\vb{x}}{t}}_{S'} + 2 \vb*{\omega} \cross \pqty{\dv{\vb{x}}{t}}_{S'} + \dot{\vb*{\omega}} \cross \vb{x} + \vb*{\omega} \cross \pqty{\vb*{\omega} \cross \vb{x}}.}
    \]
\end{proposition}

\begin{proof}
    Using Newton's Law, we have
    \begin{align*}
        \ddot{\vb{x}} &= \ddot{x}_i \vb{e}_i \\
        &= \ddot{x}_i' \vb{e}_i' + \dot{x}_i' \dot{\vb{e}}_i' + \dot{\vb*{\omega}} \cross \vb{x} + \vb*{\omega} \cross \dot{\vb{x}} \\
        &= \pqty{\dv[2]{\vb{x}}{t}}_{S'} + \dot{x}_i' \pqty{\vb*{\omega} \cross \vb{e}_i'} + \dot{\vb*{\omega}} \cross \vb{x} + \vb*{\omega} \cross \dot{\vb{x}} \\
        &= \pqty{\dv[2]{\vb{x}}{t}}_{S'} + \vb*{\omega} \cross \pqty{\dv{\vb{x}}{t}}_{S'} + \dot{\vb*{\omega}} \cross \vb{x} + \vb*{\omega} \cross \dot{\vb{x}}.
    \end{align*}

    Hence, we have
    \begin{align*}
        \pqty{\dv[2]{\vb{x}}{t}}_{S} &= \pqty{\dv[2]{\vb{x}}{t}}_{S'} + \dot{\vb*{\omega}} \cross \vb{x} + \vb*{\omega} \cross \dot{\vb{x}} + \vb*{\omega} \cross \pqty{\dv{\vb{x}}{t}}_{S'} \\
        &= \pqty{\dv[2]{\vb{x}}{t}}_{S'} + \dot{\vb*{\omega}} \cross \vb{x} + \vb*{\omega} \cross \pqty{\pqty{\dv{\vb{x}}{t}}_{S'} + \vb*{\omega} \cross \vb{x}} + \vb*{\omega} \cross \pqty{\dv{\vb{x}}{t}}_{S'} \\
        &= \pqty{\dv[2]{\vb{x}}{t}}_{S'} + 2 \vb*{\omega} \cross \pqty{\dv{\vb{x}}{t}}_{S'} + \dot{\vb*{\omega}} \cross \vb{x} + \vb*{\omega} \cross \pqty{\vb*{\omega} \cross \vb{x}}.
    \end{align*}
\end{proof}

Putting Newton's equations into this, we get \autoref{cor:fictitious-force-rotating-frame}.

\begin{corollary}[Fictitious Forces in Rotating Frame]
    \label{cor:fictitious-force-rotating-frame}%
    In a rotating frame, Newton's Law involving the fictitious forces takes the form
    \[
        \boxed{m \pqty{\dv[2]{\vb{x}}{t}}_{S'} = \vb{F} - m \dot{\vb*{\omega}} \cross \vb{x} - 2 m \vb*{\omega} \cross \pqty{\dv{\vb{x}}{t}}_{S'} - m \vb*{\omega} \cross \pqty{\vb*{\omega} \cross \vb{x}}.}
    \]
\end{corollary}

\begin{definition}[Fictitious Forces: Euler/Coriolis/Centrifugal]
    The fictitious forces involved in \autoref{cor:fictitious-force-rotating-frame} are named:
    \begin{itemize}
        \item \emph{Euler Force}: \(m \dot{\vb*{\omega}} \cross \vb{x}\);
        \item \emph{Coriolis Force}: \(2 m \vb*{\omega} \cross \pqty{\dv{\vb{x}}{t}}_{S'}\);
        \item \emph{Centrifugal Force}: \(m \vb*{\omega} \cross \pqty{\vb*{\omega} \cross \vb{x}}\).
    \end{itemize}
\end{definition}

In other words, a free particle with \(\vb{F} = \vb{0}\) does \emph{not} move on a straight line in a rotating frame.

\begin{example}[Rotation of Earth]
    \label{ex:rotation-of-earth}%
    For the rotation of Earth,
    \[
        \omega_{\text{rot}} = \frac{2 \pi}{\SI{1}{\day}} \approx \SI{7e5}{\per\second},
    \]
    and
    \[
        R_{\text{Earth}} \approx \SI{6e3}{\kilo\metre}.
    \]

    We neglect small `wobbles' of Earth, which means \(\dot{\vb*{\omega}} = \vb{0}\), so we do not experience Euler Force on Earth.
\end{example}

\subsection{Centrifugal Force}

The centrifugal force
\[
    \vb{F}_{\text{cent}} = - m \vb*{\omega} \cross \pqty{\vb*{\omega} \cross \vb{x}}
\]
points away from axis of rotation, with magnitude
\[
    \abs{\vb{F}_{\text{cent}}} = m \omega^2 r \cos \theta = m \omega^2 d.
\]

\begin{figure}[ht!]
    \centering

    \caption{Centrifugal Force}
    \label{fig:centrifugal-force}%
\end{figure}

\begin{proposition}[Centrifugal Potential]
    \label{prop:centrifugal-potential}%
    The centrifugal force is conservative, satisfying
    \[
        \vb{F}_{\text{cent}} = - \grad V_{\text{cent}},
    \]
    where the centrifugal potential \(V_{\text{cent}}\) is given by
    \[
        V_{\text{cent}} = - \frac{m}{2} \abs{\vb*{\omega} \cross \vb{x}}^2 = - \frac{m}{2} \omega^2 d^2.
    \]
\end{proposition}

The centrifugal potential energy can be lowered by moving away from the axis of rotation.

\begin{figure}[ht!]
    \centering

    \caption{Hanging String for Centrifugal Force}
    \label{fig:hanging-string-centrifugal}%
\end{figure}

\begin{example}[Hanging String]
    \label{ex:hanging-string}%
    We aim to calculate how much centrifugal force displaces a hanging string by directly pointing towards the centre of the force, as setup in \autoref{fig:hanging-string-centrifugal}.

    Forces acting on the particle (i.e. weight at the end of the string) satisfies
    \[
        \vb{g} = - m g \unitv{r}.
    \]

    As the force diagram in \autoref{fig:hanging-string-centrifugal} demonstrates, we have
    \begin{align*}
        \vb{F}_{\text{cent}} &= - m \vb*{\omega} \cross \pqty{\vb*{\omega} \cross \vb{x}} \\
        &= m \omega^2 r \cos \theta \pqty{\cos \theta \unitv{r} - \sin \theta \unitv*{\theta}}.
    \end{align*}

    To hold the string together, there must be a force exerted by the molecules in the string that balances the other forces, and macroscopically, we call it the tension \(\vb{T}\). Therefore, we have
    \[
        \vb{T} = T \cos \phi \unitv{r} + T \sin \phi \unitv*{\theta}.
    \]

    In equilibrium, balancing gives
    \[
        \vb{g} + \vb{F}_{\text{cent}} + \vb{T} = \vb{0}.
    \]

    Two equations in the \(\unitv{r}\) and \(\unitv*{\theta}\) components allows us to solve for the unknowns \(T\) and \(\phi\), and hence,
    \[
        \tan \phi = \frac{\omega^2 R \cos \theta \sin \theta}{g - \omega^2 R \cos^2 \theta}.
    \]

    Putting in familiar values, at the equator, \(\theta = 0\), gravity would be a bit ``weaker'', but \(\phi = 0\). At the North Pole, \(\theta = \frac{\pi}{2}\), and the object receives no centrifugal force, so \(\phi = 0\) indeed.

    When \(\theta \approx \SI{45}{\degree}\), we have \(\phi \approx 10^{-4}\), which is a very small effect.
\end{example}

\subsection{Coriolis Force}

The Coriolis force satisfies
\[
    \vb{F}_{\text{cor}} = - 2 m \vb*{\omega} \cross \vb{v},
\]
where \(\vb{v}\) is the velocity measured in the rotating frame (e.g., the Earth).

We note the similarity with the Lorentz force, with the magnetic flux density \(\vb{B}\) substituted by \(\vb*{\omega}\). Therefore, moving particles will turn in circles.

\begin{figure}[ht!]
    \centering

    \caption{Formation of Hurricanes}
    \label{fig:hurricanes}%
\end{figure}

\begin{example}[Formation of Hurricanes, Qualitative Discussion]
    \label{ex:hurricanes}%
    When low pressure region forms, particles would naturally move in, while Coriolis force will bend them in, as demonstrated in \autoref{fig:hurricanes}.

    In the Northern Hemisphere, each molecule of air is bent clockwise, by considering right-hand rule with \(-\vb*{\omega}\) pointing into the plane. This gives anti-clockwise swirling motion. Indeed, hurricanes in the Northern Hemisphere rotate anti-clockwise, while they rotate clockwise in the Southern Hemisphere.

    Motion along the Earth's surface is in general not perpendicular, to \(\vb*{\omega}\). The effect of Coriolis force is weaker near the equator, so hurricanes barely occur near equators.

    In general, the magnitude \(\vb*{\omega} \cross \vb{v}\) can be substantial near the equator, but it pushes the particles vertically, but it is small compared to gravity.
\end{example}

\begin{figure}[ht!]
    \centering

    \caption{Dropping a Ball from a Tower}
    \label{fig:dropping-ball-from-tower}%
\end{figure}

\begin{example}[Dropping a Ball from a Tower]
    \label{ex:dropping-ball-from-tower}%
    We drop a ball from a tower on the equator, and we study where it falls.

    First, we study the answer qualitatively, using conservation of angular momentum. Initially, the angular momentum is given by
    \[
        l = \omega \pqty{R + h}^2,
    \]
    and as the ball falls, the distance to axis decreases, so the angular velocity must speed up to conserve \(l\). For example, at the foot,
    \[
        l = \omega' R^2,
    \]
    so \(\omega' > \omega\), and hence the ball rotates faster than Earth, and must fall in front of the tower.

    Quantitatively, neglecting centrifugal force, we have
    \[
        \ddot{\vb{x}} = \vb{g} - 2 \vb*{\omega} \cross \dot{\vb{x}}
    \]
    in the rotating frame. Integrating once, we get
    \[
        \dot{\vb{x}} = \vb{g} t - 2 \vb*{\omega} \cross \pqty{\vb{x} - \vb{x}_0},
    \]
    where \(\vb{x}_0\) is the initial position (an integration constant). Substituting back into the original equation, we get
    \[
        \ddot{\vb{x}} = \vb{g} - 2 \vb*{\omega} \cross \vb{g} t + 4 \vb*{\omega} \cross \pqty{\vb*{\omega} \cross \pqty{\vb{x} - \vb{x}_0}}.
    \]

    The final term has the same order as the centrifugal force, and acts in vertical direction, so is small compared to gravity. Therefore, we have
    \[
        \ddot{\vb{x}} = \vb{g} - 2 \vb*{\omega} \cross \vb{g} t.
    \]

    We integrate twice to get
    \[
        \vb{x} = \vb{x}_0 + \frac{1}{2} \vb{g} t^2 - \frac{1}{3} \pqty{\vb*{\omega} \cross \vb{g}} t^3.
    \]

    Constructing the basis in \autoref{fig:dropping-ball-from-tower}, using \(\vb{e}_1, \vb{e}_2, \vb{e}_3\) as North, West and Up respectively, we have
    \begin{align*}
        \vb*{\omega} &= \omega \vb{e}_1, \\
        \vb{g} &= - g \vb{e}_3, \\
        \vb{x}_0 &= \pqty{R + h} \vb{e}_3.
    \end{align*}

    Hence, we may solve that
    \[
        \vb{x} = \pmqty{
            0 \\
            - \frac{1}{3} \omega g t^3 \\
            R + h - \frac{1}{2} g t^2
        }.
    \]

    When the object hits the ground, we have \(\frac{1}{2} g t^2 = h\). Clearly, the positive \(t\) solution gives \(x_2 < 0\) at this time, so it has fallen `east' to the tower, in its front.
\end{example}

\begin{figure}[ht!]
    \centering

    \caption{Foucault's Pendulum}
    \label{fig:foucault}%
\end{figure}

\begin{example}[Foucault's Pendulum]
    \label{ex:foucault}%
    Foucault's pendulum demonstrates the rotation of the Earth.

    Consider an extremal case -- if we hang a pendulum at the North Pole, as the Earth rotates under the pendulum, from the point of view of someone on the Earth, it will look like the pendulum rotates.

    At a general latitude, we build the coordinate axis similarly as in \autoref{ex:dropping-ball-from-tower}, so we have
    \[
        \vb*{\omega} = \pmqty{\omega \cos \theta \\ 0 \\ \omega \sin \theta}
    \]
    where \(\theta\) is the latitude of the object.

    The tension in the string satisfies
    \[
        \vb{T} = T \pmqty{- \frac{x}{l} \\ - \frac{y}{l} \\ \frac{l - z}{l}},
    \]
    with the geometrical constraint
    \[
        x^2 + y^2 + \pqty{l - z}^2 = l^2.
    \]

    With the equations of motion (giving 3 ODEs), in addition to the constraints, it suffices to solve the four unknowns, \(x, y, z\) and \(T\). The general strategy is to:
    \begin{enumerate}
        \item solve constraints of this equation, giving \(z = z \pqty{x, y}\);
        \item put into the ODEs;
        \item eliminate \(T\);
        \item form 2 ODEs for \(x\pqty{t}\) and \(y\pqty{t}\).
    \end{enumerate}

    The solution is tedious, but the upshot is that the pendulum, in general, follows an ellipse in the \(\pqty{x, y}\) plane that slowly rotates, i.e.,
    \[
        x + iy = \exp \pqty{- i \omega t \sin \theta} \cdot \bqty{a \cos \pqty{\sqrt{\frac{g}{l}} t} + b \sin \pqty{\sqrt{\frac{g}{l}}} t},
    \]
    where \(a\) and \(b\) are complex constants, and \(\exp\) gives the rotation of the ellipse. This means, the period of rotation is
    \[
        \frac{\SI{24}{\hour}}{\sin \theta}.
    \]

    The extremal cases \(\theta = \frac{\pi}{2}\) and \(\theta \to 0^+\) indeed satisfies this equation intuitively.
\end{example}